\documentclass[../BTL.tex]{subfiles}
\begin{document}
\section{Phân tích hệ thống}
\subsection{Xác định các tác nhân trong hệ thống}

Hệ thống chấm bài tự động có bốn tác nhân, trong đó User là tác nhân trừu tượng còn Quản trị viên, Giảng viên và Sinh viên là các tác nhân cụ thể kế thừa từ User:
\begin{itemize}
    \item Người dùng (User): tác nhân trừu tượng đại diện cho mọi tài khoản truy cập hệ thống; các tác nhân còn lại đều kế thừa User nên đều thực hiện được hai usecase chung là đăng nhập và đổi mật khẩu.
    \item Quản trị viên (Admin): quản trị tài khoản ở mức hệ thống; là tác nhân duy nhất được phép nhập danh sách tài khoản người dùng từ tệp CSV để cấp tài khoản ban đầu cho giảng viên và sinh viên.
    \item Giảng viên (Lecturer): quản lý lớp học của mình gồm tạo lớp, nhập danh sách sinh viên, cấu hình thành phần điểm và xem bảng điểm; quản lý bài tập kèm các bước chấm tự động; xem kết quả chấm bài của cả lớp.
    \item Sinh viên (Student): tham gia các lớp đã được ghi danh, nộp bài thực hành để hệ thống chấm điểm tự động và xem điểm, kết quả chấm của bản thân.
\end{itemize}

\subsection{Biểu đồ ca chức năng (usecase) tổng quan của hệ thống}

Hình \ref{fig:usecase-tong-quan} mô tả các ca chức năng tổng quan của hệ thống và các tác nhân tương tác với chúng.

\begin{figure}[ht]
\centering
\includegraphics[width=\textwidth]{Hinhve/Hinh2_6_usecase_tong_quan}
\caption{Biểu đồ usecase tổng quan của hệ thống.}
\label{fig:usecase-tong-quan}
\end{figure}

Mô tả các usecase của hệ thống:
\begin{itemize}
    \item Đăng nhập: Usecase này cho phép người dùng (quản trị viên, giảng viên và sinh viên) truy cập vào hệ thống bằng tài khoản đã có để thực hiện được các chức năng theo đúng quyền quy định của mình.
    \item Đổi mật khẩu: Usecase này cho phép người dùng đã đăng nhập thay đổi mật khẩu tài khoản của mình để tiếp tục sử dụng hệ thống một cách an toàn với thông tin đăng nhập mới.
    \item Import users accounts (nhập danh sách tài khoản): Usecase này cho phép quản trị viên tạo hàng loạt tài khoản cho giảng viên và sinh viên từ tệp CSV để cấp tài khoản ban đầu mà không phải tạo thủ công từng tài khoản. Mỗi tài khoản mới nhận một mật khẩu tạm và bắt buộc đổi mật khẩu ở lần đăng nhập đầu tiên.
    \item Classes management (quản lý lớp học): Usecase này cho phép giảng viên tạo và lưu trữ lớp học của mình, nhập danh sách sinh viên vào lớp, cấu hình các thành phần điểm và xem bảng điểm của lớp để quản lý toàn bộ hoạt động của một lớp học trong một học kỳ.
    \item Assignment management (quản lý bài tập): Usecase này cho phép giảng viên tạo, sửa, công bố và xoá bài tập thực hành trong lớp của mình, kèm theo các bước chấm tự động để hệ thống chấm bài nộp của sinh viên theo đúng yêu cầu đề ra.
    \item Review class assignment results (xem kết quả chấm của cả lớp): Usecase này cho phép giảng viên xem kết quả chấm bài tập của toàn bộ sinh viên trong lớp, bao gồm điểm tổng hợp của từng sinh viên, điểm từng bài nộp và chi tiết từng bước chấm, để theo dõi và đánh giá tình hình học tập của lớp.
    \item Review assignment results (xem chi tiết kết quả chấm một bài tập): Usecase này là bước chung khi xem chi tiết kết quả chấm của một bài tập: giảng viên đi tới usecase này từ Review class assignment results, còn sinh viên đi tới từ Review own assignment results; phạm vi kết quả nhìn thấy phụ thuộc vào tác nhân gọi nó.
    \item Submit assignments (nộp bài): Usecase này cho phép sinh viên nộp bài thực hành của mình dưới dạng tệp nén để hệ thống tải lên và chấm điểm tự động, đồng thời theo dõi trạng thái chấm và xem lại kết quả của các lần nộp.
    \item View personal scores and results (xem điểm tổng kết cá nhân): Usecase này cho phép sinh viên xem điểm thành phần, điểm tổng kết, điểm chữ và điểm hệ 4 của mình trong mỗi lớp đã ghi danh để tự đánh giá kết quả học tập của bản thân.
    \item Review own assignment results (xem kết quả chấm bài tập của bản thân): Usecase này cho phép sinh viên xem điểm và chi tiết từng bước chấm của các bài tập mình đã nộp, giúp hiểu rõ vì sao bài làm đạt hay chưa đạt yêu cầu của đề bài.
\end{itemize}

\subsection{Kịch bản chuẩn các chức năng}

Sau khi xác định các tác nhân và biểu đồ ca chức năng tổng quan, mục này đặc tả kịch bản chuẩn của từng chức năng trong hệ thống. Các chức năng được nhóm theo tác nhân sử dụng: nhóm chức năng chung cho mọi người dùng, nhóm chức năng của quản trị viên, nhóm chức năng của giảng viên và nhóm chức năng của sinh viên. Mỗi chức năng gồm biểu đồ ca chức năng chi tiết (nếu có), danh sách các ca chức năng con và bảng kịch bản chuẩn; bảng kịch bản mô tả tác nhân, điều kiện tiên quyết, luồng sự kiện chính, hậu điều kiện và các ngoại lệ có thể xảy ra khi thực hiện chức năng. Usecase Review assignment results là bước chung khi xem chi tiết một bài nộp, nên được mô tả ngay trong kịch bản Review class assignment results của giảng viên và Review own assignment results của sinh viên thay vì lập bảng riêng.

\subsubsection{Kịch bản các chức năng chung}

\paragraph{Chức năng Đăng nhập.}
Bảng \ref{tab:uc-dang-nhap} mô tả kịch bản chuẩn của usecase Đăng nhập.

\begin{table}[!ht]
\centering
\begin{tabularx}{\textwidth}{|p{3.2cm}|X|}
\hline
Tên mục & Nội dung mô tả \\ \hline
Tên kịch bản & Đăng nhập \\ \hline
Tác nhân & Người dùng (quản trị viên, giảng viên và sinh viên) \\ \hline
Tiền điều kiện & Người dùng đã được cấp tài khoản trên hệ thống (do quản trị viên tạo hoặc nhập từ tệp CSV). \\ \hline
Kịch bản chính &
\begin{enumerate}[leftmargin=*, nosep]
    \item Người dùng mở hệ thống và chọn đăng nhập; hệ thống hiển thị trang đăng nhập.
    \item Người dùng nhập tên đăng nhập và mật khẩu; hệ thống kiểm tra thông tin đăng nhập.
    \item Nếu tài khoản còn yêu cầu đổi mật khẩu bắt buộc (tài khoản cấp mới bằng mật khẩu tạm), hệ thống yêu cầu người dùng đổi mật khẩu trước khi tiếp tục.
    \item Hệ thống xác nhận đăng nhập thành công và thiết lập phiên làm việc cho người dùng.
    \item Ứng dụng mở giao diện tương ứng với vai trò của người dùng (quản trị viên, giảng viên hoặc sinh viên).
\end{enumerate} \\ \hline
Hậu điều kiện & Người dùng có phiên làm việc hợp lệ; các chức năng tiếp theo được phân quyền theo vai trò. \\ \hline
Kịch bản ngoại lệ &
\begin{itemize}[leftmargin=*, nosep]
    \item Tên đăng nhập hoặc mật khẩu không đúng: trang đăng nhập báo lỗi và cho phép nhập lại.
    \item Nhập sai mật khẩu quá số lần quy định: tài khoản bị tạm khoá để chống dò mật khẩu.
    \item Phiên làm việc hết hạn: hệ thống yêu cầu người dùng đăng nhập lại.
\end{itemize} \\ \hline
\end{tabularx}
\caption{Kịch bản chức năng Đăng nhập.}
\label{tab:uc-dang-nhap}
\end{table}

\paragraph{Chức năng Đổi mật khẩu.}
Bảng \ref{tab:uc-doi-mat-khau} mô tả kịch bản chuẩn của usecase Đổi mật khẩu.

\begin{table}[!ht]
\centering
\begin{tabularx}{\textwidth}{|p{3.2cm}|X|}
\hline
Tên mục & Nội dung mô tả \\ \hline
Tên kịch bản & Đổi mật khẩu \\ \hline
Tác nhân & Người dùng đã đăng nhập (quản trị viên, giảng viên và sinh viên) \\ \hline
Tiền điều kiện & Người dùng đang có phiên làm việc hợp lệ. \\ \hline
Kịch bản chính &
\begin{enumerate}[leftmargin=*, nosep]
    \item Người dùng mở menu tài khoản và chọn chức năng Đổi mật khẩu.
    \item Hệ thống hiển thị hộp thoại gồm ba trường: mật khẩu hiện tại, mật khẩu mới và xác nhận mật khẩu mới.
    \item Người dùng nhập thông tin; giao diện kiểm tra ngay độ dài tối thiểu và sự trùng khớp giữa mật khẩu mới với phần xác nhận trước khi gửi yêu cầu.
    \item Hệ thống xác minh mật khẩu hiện tại, sau đó đặt mật khẩu mới cho tài khoản.
    \item Hệ thống thông báo đổi mật khẩu thành công và đóng hộp thoại.
\end{enumerate} \\ \hline
Hậu điều kiện & Mật khẩu mới được áp dụng cho lần đăng nhập sau; phiên làm việc hiện tại không bị kết thúc nên người dùng tiếp tục làm việc bình thường. \\ \hline
Kịch bản ngoại lệ &
\begin{itemize}[leftmargin=*, nosep]
    \item Mật khẩu hiện tại không đúng: hệ thống báo lỗi và yêu cầu nhập lại.
    \item Mật khẩu mới quá ngắn hoặc xác nhận không khớp: giao diện chặn và không gửi yêu cầu.
    \item Hệ thống tạm thời không xử lý được yêu cầu: báo lỗi và yêu cầu người dùng thử lại sau.
\end{itemize} \\ \hline
\end{tabularx}
\caption{Kịch bản chức năng Đổi mật khẩu.}
\label{tab:uc-doi-mat-khau}
\end{table}

\subsubsection{Kịch bản các chức năng của quản trị viên}

\paragraph{Chức năng Import users accounts.}
Bảng \ref{tab:uc-import-users} mô tả kịch bản chuẩn của usecase Import users accounts.

\begin{table}[!ht]
\centering
\begin{tabularx}{\textwidth}{|p{3.2cm}|X|}
\hline
Tên mục & Nội dung mô tả \\ \hline
Tên kịch bản & Import users accounts (nhập danh sách tài khoản) \\ \hline
Tác nhân & Quản trị viên \\ \hline
Tiền điều kiện & Quản trị viên đã đăng nhập bằng tài khoản có vai trò quản trị. \\ \hline
Kịch bản chính &
\begin{enumerate}[leftmargin=*, nosep]
    \item Quản trị viên mở chức năng nhập danh sách tài khoản và chọn tệp CSV gồm các cột: tên đăng nhập, họ tên, email và vai trò (giảng viên hoặc sinh viên; bỏ trống mặc định là sinh viên).
    \item Hệ thống kiểm tra giới hạn kích thước tệp và số dòng, sau đó xử lý tuần tự từng dòng.
    \item Với mỗi dòng hợp lệ, hệ thống kiểm tra trùng tên đăng nhập và trùng email, tạo tài khoản mới với mật khẩu tạm kèm yêu cầu đổi mật khẩu ở lần đăng nhập đầu tiên, rồi gán vai trò tương ứng.
    \item Hệ thống trả về báo cáo gồm số tài khoản đã tạo theo từng vai trò, số dòng bị bỏ qua và danh sách dòng lỗi kèm lý do.
\end{enumerate} \\ \hline
Hậu điều kiện & Các tài khoản mới sẵn sàng đăng nhập bằng mật khẩu tạm; người dùng bắt buộc đổi mật khẩu ở lần đăng nhập đầu tiên. \\ \hline
Kịch bản ngoại lệ &
\begin{itemize}[leftmargin=*, nosep]
    \item Tệp vượt quá giới hạn kích thước hoặc số dòng: hệ thống từ chối cả tệp và báo lỗi.
    \item Dòng thiếu tên đăng nhập hoặc email, hoặc vai trò không hợp lệ: dòng đó bị bỏ qua và được ghi vào báo cáo kèm lý do, các dòng khác vẫn được xử lý.
    \item Tên đăng nhập hoặc email đã tồn tại: dòng đó bị bỏ qua; nhập lại cùng tệp vẫn an toàn.
\end{itemize} \\ \hline
\end{tabularx}
\caption{Kịch bản chức năng Import users accounts.}
\label{tab:uc-import-users}
\end{table}

\subsubsection{Kịch bản các chức năng của giảng viên}

\paragraph{Chức năng Classes management.}

Hình \ref{fig:uc-classes-management} mô tả chi tiết usecase Classes management của giảng viên.

\begin{figure}[ht]
\centering
\includegraphics[width=\textwidth]{Hinhve/Hinh2_7_usecase_quan_ly_lop_hoc}
\caption{Biểu đồ usecase chi tiết Classes management của giảng viên.}
\label{fig:uc-classes-management}
\end{figure}

Usecase này được phân rã thành các usecase con sau:
\begin{itemize}
    \item Lecturer login (đăng nhập của giảng viên): usecase con được bao gồm (quan hệ include) trong Classes management, đảm bảo giảng viên đã đăng nhập đúng vai trò trước khi quản lý lớp học; Lecturer login kế thừa (quan hệ generalization) từ usecase Đăng nhập.
    \item View class's students list (xem danh sách sinh viên của lớp): mở rộng (quan hệ extend) từ Classes management khi giảng viên cần xem danh sách sinh viên hiện có trong lớp.
    \item Import students to class (nhập sinh viên vào lớp): mở rộng từ Classes management khi giảng viên thêm danh sách sinh viên vào lớp từ tệp CSV.
    \item Config course score components (cấu hình thành phần điểm): mở rộng từ Classes management khi giảng viên khai báo các thành phần điểm và trọng số tương ứng của lớp.
    \item View class transcript (xem bảng điểm của lớp): bao gồm (quan hệ include) Config course score components, cho phép giảng viên xem điểm tổng hợp của toàn bộ sinh viên theo đúng cấu trúc điểm đã cấu hình.
\end{itemize}

Kịch bản chuẩn của chức năng được trình bày trong Bảng \ref{tab:uc-classes-management}.

\begin{table}[!ht]
\centering
\begin{tabularx}{\textwidth}{|p{3.2cm}|X|}
\hline
Tên mục & Nội dung mô tả \\ \hline
Tên kịch bản & Classes management (quản lý lớp học) \\ \hline
Tác nhân & Giảng viên \\ \hline
Tiền điều kiện & Giảng viên đã đăng nhập bằng tài khoản có vai trò giảng viên. \\ \hline
Kịch bản chính &
\begin{enumerate}[leftmargin=*, nosep]
    \item Giảng viên mở mục quản lý lớp học; hệ thống hiển thị danh sách lớp của giảng viên kèm trạng thái, có tìm kiếm theo tên lớp, học kỳ và bộ lọc theo trạng thái.
    \item Giảng viên tạo lớp mới bằng tên lớp và học kỳ; hệ thống tạo lớp ở trạng thái đang hoạt động.
    \item Giảng viên nhập danh sách sinh viên của lớp từ tệp CSV; hệ thống báo số sinh viên được thêm và số dòng bị bỏ qua.
    \item Giảng viên cấu hình các thành phần điểm của lớp (chuyên cần, bài tập, thi cuối kỳ) cùng trọng số tương ứng; hệ thống kiểm tra các ràng buộc về trọng số trước khi lưu.
    \item Giảng viên nhập điểm thủ công cho từng sinh viên; thành phần điểm bài tập được lấy tự động từ kết quả chấm bài nộp, sau đó hệ thống tính điểm tổng kết, điểm chữ và điểm hệ 4.
    \item Giảng viên xem bảng điểm của cả lớp hoặc lưu trữ lớp khi kết thúc học kỳ.
\end{enumerate} \\ \hline
Hậu điều kiện & Lớp học, danh sách sinh viên, cấu trúc điểm và điểm số được lưu trong hệ thống; lớp đã lưu trữ chỉ còn xem và không thể mở lại. \\ \hline
Kịch bản ngoại lệ &
\begin{itemize}[leftmargin=*, nosep]
    \item Tên lớp trùng trong cùng học kỳ: hệ thống báo lỗi và không tạo lớp.
    \item Tệp CSV sai định dạng: hệ thống báo lỗi.
    \item Cấu trúc điểm thiếu thành phần thi cuối kỳ, trọng số thành phần này thấp hơn mức tối thiểu hoặc tổng trọng số không bằng 1: hệ thống từ chối lưu.
    \item Sinh viên không thuộc danh sách lớp: hệ thống báo lỗi khi nhập điểm.
    \item Một thành phần điểm chưa có điểm: điểm tổng kết, điểm chữ và điểm hệ 4 để trống trên bảng điểm.
\end{itemize} \\ \hline
\end{tabularx}
\caption{Kịch bản chức năng Classes management.}
\label{tab:uc-classes-management}
\end{table}

\paragraph{Chức năng Assignment management.}

Hình \ref{fig:uc-assignment-management} mô tả chi tiết usecase Assignment management của giảng viên.

\begin{figure}[ht]
\centering
\includegraphics[width=\textwidth]{Hinhve/Hinh2_8_usecase_quan_ly_bai_tap}
\caption{Biểu đồ usecase chi tiết Assignment management của giảng viên.}
\label{fig:uc-assignment-management}
\end{figure}

Usecase này được phân rã thành các usecase con sau:
\begin{itemize}
    \item Lecturer login (đăng nhập của giảng viên): usecase con được bao gồm (quan hệ include) trong Assignment management, đảm bảo giảng viên đã đăng nhập đúng vai trò trước khi quản lý bài tập.
    \item View class (chọn lớp): bao gồm (quan hệ include), cho phép giảng viên chọn lớp cần quản lý bài tập.
    \item author assignment exercises (soạn nội dung bài tập): mở rộng từ Assignment management khi giảng viên soạn đề bài và các bước chấm tự động cho bài tập.
    \item View assignment grading results (xem kết quả chấm bài tập): mở rộng từ Assignment management khi giảng viên muốn xem kết quả chấm của bài tập.
\end{itemize}

Kịch bản chuẩn của chức năng được trình bày trong Bảng \ref{tab:uc-assignment-management}.

\begin{table}[!ht]
\centering
\begin{tabularx}{\textwidth}{|p{3.2cm}|X|}
\hline
Tên mục & Nội dung mô tả \\ \hline
Tên kịch bản & Assignment management (quản lý bài tập) \\ \hline
Tác nhân & Giảng viên \\ \hline
Tiền điều kiện & Giảng viên đã đăng nhập; lớp được chọn thuộc giảng viên. \\ \hline
Kịch bản chính &
\begin{enumerate}[leftmargin=*, nosep]
    \item Giảng viên chọn lớp cần quản lý; hệ thống hiển thị danh sách bài tập của lớp kèm bộ lọc theo trạng thái công bố và tìm kiếm theo tiêu đề.
    \item Giảng viên tạo bài tập với tiêu đề và cấu hình chấm (chiến lược chấm, thời gian chờ và giới hạn tài nguyên khi chấm).
    \item Giảng viên soạn các bước chấm tự động cho bài tập: gửi yêu cầu kiểm thử tới ứng dụng của sinh viên, kiểm tra dữ liệu và trích xuất biến để dùng cho các bước sau; mỗi bộ bước như vậy là một kế hoạch chấm.
    \item Giảng viên chọn môi trường chạy cần thiết cho bài tập.
    \item Giảng viên công bố bài tập để sinh viên nhìn thấy và nộp được bài.
    \item Giảng viên xem kết quả chấm, sửa thông tin hoặc xoá bài tập khi không còn sử dụng.
\end{enumerate} \\ \hline
Hậu điều kiện & Bài tập cùng các bước chấm được lưu; bài tập đã công bố xuất hiện trong danh sách bài tập của sinh viên đã ghi danh. \\ \hline
Kịch bản ngoại lệ &
\begin{itemize}[leftmargin=*, nosep]
    \item Tiêu đề bài tập trùng trong cùng lớp: hệ thống báo lỗi và không tạo bài tập.
    \item Lớp không thuộc giảng viên: hệ thống báo không tìm thấy.
    \item Cấu hình chấm thiếu thông tin bắt buộc hoặc bước chấm trùng thứ tự: hệ thống báo lỗi kèm vị trí vi phạm.
    \item Bài tập đã xoá: không còn xuất hiện với cả giảng viên và sinh viên.
\end{itemize} \\ \hline
\end{tabularx}
\caption{Kịch bản chức năng Assignment management.}
\label{tab:uc-assignment-management}
\end{table}

\paragraph{Chức năng Review class assignment results.}

Hình \ref{fig:uc-review-class-results} mô tả chi tiết usecase Review class assignment results của giảng viên.

\begin{figure}[ht]
\centering
\includegraphics[width=\textwidth]{Hinhve/Hinh2_9_usecase_xem_ket_qua_lop}
\caption{Biểu đồ usecase chi tiết Review class assignment results của giảng viên.}
\label{fig:uc-review-class-results}
\end{figure}

Usecase này được phân rã thành các usecase con sau:
\begin{itemize}
    \item Lecturer login (đăng nhập của giảng viên): usecase con được bao gồm (quan hệ include), đảm bảo giảng viên đã đăng nhập đúng vai trò trước khi xem kết quả chấm của lớp.
    \item View classes (chọn lớp): bao gồm (quan hệ include), cho phép giảng viên chọn lớp cần xem kết quả chấm.
    \item View assignments (chọn bài tập): bao gồm (quan hệ include), cho phép giảng viên chọn bài tập cần xem kết quả chấm.
    \item Review detail assignment step results (xem chi tiết kết quả từng bước chấm): mở rộng từ Review class assignment results khi giảng viên cần xem chi tiết kết quả từng bước chấm của bài nộp.
\end{itemize}

Kịch bản chuẩn của chức năng được trình bày trong Bảng \ref{tab:uc-review-class-results}.

\begin{table}[!ht]
\centering
\begin{tabularx}{\textwidth}{|p{3.2cm}|X|}
\hline
Tên mục & Nội dung mô tả \\ \hline
Tên kịch bản & Review class assignment results (xem kết quả chấm của cả lớp) \\ \hline
Tác nhân & Giảng viên \\ \hline
Tiền điều kiện & Giảng viên sở hữu lớp; bài tập đã có ít nhất một bài nộp được chấm. \\ \hline
Kịch bản chính &
\begin{enumerate}[leftmargin=*, nosep]
    \item Giảng viên chọn lớp và bài tập cần xem kết quả.
    \item Hệ thống hiển thị điểm tổng hợp của từng sinh viên trong lớp cùng điểm từng bài nộp theo từng kế hoạch chấm.
    \item Giảng viên lọc theo mã sinh viên khi cần xem kết quả của một sinh viên cụ thể.
    \item Giảng viên mở chi tiết một bài nộp để xem kết quả từng bước chấm, gồm yêu cầu đã gửi, kết quả nhận về và nhận xét của từng bước.
    \item Giảng viên xem danh sách toàn bộ bài nộp của bài tập.
\end{enumerate} \\ \hline
Hậu điều kiện & Không thay đổi dữ liệu; kết quả chỉ dùng để theo dõi và đánh giá tình hình học tập của lớp. \\ \hline
Kịch bản ngoại lệ &
\begin{itemize}[leftmargin=*, nosep]
    \item Bài tập không thuộc giảng viên: hệ thống báo không tìm thấy.
    \item Mã sinh viên không có trong lớp: hệ thống trả về danh sách rỗng.
    \item Bài tập chưa có bài nộp nào: danh sách kết quả rỗng.
\end{itemize} \\ \hline
\end{tabularx}
\caption{Kịch bản chức năng Review class assignment results.}
\label{tab:uc-review-class-results}
\end{table}

\subsubsection{Kịch bản các chức năng của sinh viên}

\paragraph{Chức năng Submit assignments.}

Hình \ref{fig:uc-submit-assignment} mô tả chi tiết usecase Submit assignments của sinh viên.

\begin{figure}[ht]
\centering
\includegraphics[width=\textwidth]{Hinhve/Hinh2_10_usecase_nop_bai}
\caption{Biểu đồ usecase chi tiết Submit assignments của sinh viên.}
\label{fig:uc-submit-assignment}
\end{figure}

Usecase này được phân rã thành các usecase con sau:
\begin{itemize}
    \item Student login (đăng nhập của sinh viên): usecase con được bao gồm (quan hệ include) trong Submit assignments, đảm bảo sinh viên đã đăng nhập đúng vai trò trước khi nộp bài; Student login kế thừa (quan hệ generalization) từ usecase Đăng nhập.
    \item View class (chọn lớp): bao gồm (quan hệ include), cho phép sinh viên chọn lớp có bài tập cần nộp.
    \item View assignments (chọn bài tập): bao gồm (quan hệ include), cho phép sinh viên chọn bài tập cần nộp.
    \item View assignment grading results (xem kết quả chấm bài tập): mở rộng từ Submit assignments khi sinh viên muốn xem lại kết quả chấm của bài đã nộp.
\end{itemize}

Kịch bản chuẩn của chức năng được trình bày trong Bảng \ref{tab:uc-submit-assignments}.

\begin{table}[!ht]
\centering
\begin{tabularx}{\textwidth}{|p{3.2cm}|X|}
\hline
Tên mục & Nội dung mô tả \\ \hline
Tên kịch bản & Submit assignments (nộp bài) \\ \hline
Tác nhân & Sinh viên \\ \hline
Tiền điều kiện & Sinh viên đã đăng nhập; đã được ghi danh vào lớp; bài tập đã được công bố. \\ \hline
Kịch bản chính &
\begin{enumerate}[leftmargin=*, nosep]
    \item Sinh viên chọn lớp và bài tập cần nộp; hệ thống hiển thị đề bài dưới dạng các bước yêu cầu, đã ẩn phần cấu hình chấm nội bộ của giảng viên.
    \item Sinh viên chọn kế hoạch chấm cần nộp, hoặc nộp cho toàn bộ bài tập, và chuẩn bị tệp nén bài làm.
    \item Hệ thống cấp đường dẫn tải lên an toàn cho bài nộp.
    \item Sinh viên tải tệp nén bài làm lên hệ thống; việc tải lên hoàn tất sẽ tự động kích hoạt chấm bài, sinh viên không cần xác nhận thêm.
    \item Hệ thống chạy chấm trong môi trường cách ly theo các bước đã soạn: khởi động ứng dụng của sinh viên, gửi yêu cầu kiểm thử và kiểm tra dữ liệu, sau đó tổng hợp điểm.
    \item Sinh viên theo dõi trạng thái chấm và xem điểm, chi tiết từng bước chấm khi hoàn tất.
\end{enumerate} \\ \hline
Hậu điều kiện & Bài nộp được lưu cùng kết quả chấm cho kế hoạch đã nộp; nếu chấm thất bại, bài nộp giữ trạng thái thất bại kèm nhật ký để sinh viên xem lại và nộp lại. \\ \hline
Kịch bản ngoại lệ &
\begin{itemize}[leftmargin=*, nosep]
    \item Sinh viên chưa ghi danh hoặc bài tập chưa công bố: hệ thống báo không tìm thấy.
    \item Tệp nộp sai định dạng hoặc tải lên thất bại: hệ thống báo lỗi và sinh viên thực hiện lại.
    \item Quá trình chấm lỗi (ứng dụng không khởi động được, hết thời gian chạy hoặc bước chấm sai cấu hình): bài nộp chuyển trạng thái thất bại kèm nhật ký lỗi.
\end{itemize} \\ \hline
\end{tabularx}
\caption{Kịch bản chức năng Submit assignments.}
\label{tab:uc-submit-assignments}
\end{table}

\paragraph{Chức năng View personal scores and results.}

Hình \ref{fig:uc-view-scores} mô tả chi tiết usecase View personal scores and results của sinh viên.

\begin{figure}[ht]
\centering
\includegraphics[width=\textwidth]{Hinhve/Hinh2_11_usecase_xem_diem_ca_nhan}
\caption{Biểu đồ usecase chi tiết View personal scores and results của sinh viên.}
\label{fig:uc-view-scores}
\end{figure}

Usecase này được phân rã thành các usecase con sau:
\begin{itemize}
    \item Student login (đăng nhập của sinh viên): usecase con được bao gồm (quan hệ include), đảm bảo sinh viên đã đăng nhập đúng vai trò trước khi xem điểm cá nhân.
    \item View classes (chọn lớp): bao gồm (quan hệ include), cho phép sinh viên chọn lớp cần xem điểm.
\end{itemize}

Kịch bản chuẩn của chức năng được trình bày trong Bảng \ref{tab:uc-view-scores}.

\begin{table}[!ht]
\centering
\begin{tabularx}{\textwidth}{|p{3.2cm}|X|}
\hline
Tên mục & Nội dung mô tả \\ \hline
Tên kịch bản & View personal scores and results (xem điểm tổng kết cá nhân) \\ \hline
Tác nhân & Sinh viên \\ \hline
Tiền điều kiện & Sinh viên đã đăng nhập và được ghi danh trong lớp. \\ \hline
Kịch bản chính &
\begin{enumerate}[leftmargin=*, nosep]
    \item Sinh viên chọn lớp đã ghi danh; hệ thống hiển thị điểm từng thành phần kèm trọng số.
    \item Hệ thống tính điểm tổng kết, điểm chữ và điểm hệ 4 khi sinh viên đã có đủ điểm các thành phần.
    \item Sinh viên xem danh sách sinh viên cùng lớp, chỉ gồm mã sinh viên và họ tên.
\end{enumerate} \\ \hline
Hậu điều kiện & Không thay đổi dữ liệu. \\ \hline
Kịch bản ngoại lệ &
\begin{itemize}[leftmargin=*, nosep]
    \item Sinh viên không ghi danh trong lớp: hệ thống báo không tìm thấy lớp.
    \item Còn thiếu điểm thành phần: điểm tổng kết, điểm chữ và điểm hệ 4 để trống; sinh viên xem lại sau khi giảng viên nhập đủ điểm.
    \item Một thành phần điểm bị điểm 0: điểm chữ xếp F bất kể điểm tổng kết.
\end{itemize} \\ \hline
\end{tabularx}
\caption{Kịch bản chức năng View personal scores and results.}
\label{tab:uc-view-scores}
\end{table}

\paragraph{Chức năng Review own assignment results.}

Hình \ref{fig:uc-review-own-results} mô tả chi tiết usecase Review own assignment results của sinh viên.

\begin{figure}[ht]
\centering
\includegraphics[width=\textwidth]{Hinhve/Hinh2_12_usecase_xem_ket_qua_bai_nop}
\caption{Biểu đồ usecase chi tiết Review own assignment results của sinh viên.}
\label{fig:uc-review-own-results}
\end{figure}

Usecase này được phân rã thành các usecase con sau:
\begin{itemize}
    \item Student login (đăng nhập của sinh viên): usecase con được bao gồm (quan hệ include), đảm bảo sinh viên đã đăng nhập đúng vai trò trước khi xem kết quả bài nộp.
    \item View classes (chọn lớp): bao gồm (quan hệ include), cho phép sinh viên chọn lớp cần xem.
    \item View assignments (chọn bài tập đã nộp): bao gồm (quan hệ include), cho phép sinh viên chọn bài tập đã nộp cần xem kết quả.
    \item View detail assignment steps result (xem chi tiết kết quả từng bước chấm): mở rộng từ Review own assignment results khi sinh viên cần xem chi tiết kết quả từng bước chấm của bài nộp.
\end{itemize}

Kịch bản chuẩn của chức năng được trình bày trong Bảng \ref{tab:uc-review-own-results}.

\begin{table}[!ht]
\centering
\begin{tabularx}{\textwidth}{|p{3.2cm}|X|}
\hline
Tên mục & Nội dung mô tả \\ \hline
Tên kịch bản & Review own assignment results (xem kết quả chấm bài tập của bản thân) \\ \hline
Tác nhân & Sinh viên \\ \hline
Tiền điều kiện & Sinh viên đã đăng nhập; bài nộp thuộc sinh viên và đã được chấm. \\ \hline
Kịch bản chính &
\begin{enumerate}[leftmargin=*, nosep]
    \item Sinh viên chọn lớp, chọn bài tập rồi chọn bài nộp của mình.
    \item Hệ thống hiển thị điểm, trạng thái chấm, nhật ký chấm và kết quả từng bước.
    \item Sinh viên mở chi tiết từng bước chấm để xem yêu cầu đã gửi, kết quả nhận về và nhận xét của bước đó; các thông tin nội bộ nhạy cảm của quá trình chấm đã được ẩn.
\end{enumerate} \\ \hline
Hậu điều kiện & Không thay đổi dữ liệu. \\ \hline
Kịch bản ngoại lệ &
\begin{itemize}[leftmargin=*, nosep]
    \item Bài nộp không thuộc sinh viên: hệ thống báo không tìm thấy.
    \item Bài nộp chưa chấm xong: hệ thống hiển thị trạng thái đang chấm; sinh viên xem lại sau.
\end{itemize} \\ \hline
\end{tabularx}
\caption{Kịch bản chức năng Review own assignment results.}
\label{tab:uc-review-own-results}
\end{table}


\end{document}
